\pagestyle{empty}
\tight
\begin{tblr}{width=\textwidth,colspec={|Q[c,m,1.9cm]|X[l,m]|},hlines,vlines,hspan=minimal,row{1}={bg=headblue},rowsep=1pt,colsep=3pt}
\SetCell{c,m}\hfil\logo\hfil & \textbf{POLITEKNIK NEGERI MALANG}\par\textbf{JURUSAN TEKNOLOGI INFORMASI}\par\textbf{PROGRAM STUDI: D4 TEKNIK INFORMATIKA} \\
\SetCell[c=2]{c} \textbf{RENCANA PEMBELAJARAN SEMESTER (RPS)} & \\
\end{tblr}
\begin{longtblr}{width=\textwidth,colspec={|Q[l,m,1.9cm]|X[0.85,l,m]|X[1.45,l,m]|X[0.75,l,m]|X[1.15,l,m]|},hlines,vlines,hspan=minimal,rowsep=1pt,colsep=2pt,row{1,3}={bg=headgray,font=\bfseries}}
MATA KULIAH & KODE & BOBOT (SKS)/jam & SEMESTER & TGL. PENYUSUNAN \\
Statistik Komputasi & RTI254008 & 2 SKS/4 jam & 4 & 1 Juli 2025 \\
OTORISASI & Dosen Pengembang RPS & Koordinator RMK & \SetCell[c=2]{l} Ka PRODI &  \\
 & Vivi Nur Wijayaningrum, S.Kom., M.Kom. & Vivi Nur Wijayaningrum, S.Kom., M.Kom. & \SetCell[c=2]{l}{\vspace{8mm}\par Dr. Ely Setyo Astuti, S.T., M.T.} &  \\
\SetCell[r=9]{l} Capaian Pembelajaran (CP) & \SetCell[c=4]{l,bg=headgray,font=\bfseries} Capaian Pembelajaran Lulusan Program Studi (CPL-Prodi) &  &  &  \\
 & CPL07 & \SetCell[c=3]{l} Mampu mengembangkan sistem komputasi cerdas dan melakukan analisis data berbasis kecerdasan buatan untuk pengambilan keputusan. &  &  \\
 & CPL09 & \SetCell[c=3]{l} Mampu menjelaskan cara kerja sistem komputer dan menerapkan pendekatan komputasi untuk menyelesaikan masalah. &  &  \\
 & \SetCell[c=4]{l,bg=headgray,font=\bfseries} Capaian Pembelajaran mata kuliah (CPL-MK) &  &  &  \\
 & CPMK0707 & \SetCell[c=3]{l} Mahasiswa mampu melakukan pemodelan dan simulasi statistika komputasi untuk analisis dan pengambilan keputusan berbasis data. &  &  \\
 & CPMK0902 & \SetCell[c=3]{l} Mahasiswa mampu menerapkan prinsip matematis dan statistik untuk analisis komputasional masalah TI. &  &  \\
 & \SetCell[c=4]{l,bg=headgray,font=\bfseries} Kemampuan akhir tiap tahapan belajar (Sub-CPMK) &  &  &  \\
 & SCPMK0707-03501 & \SetCell[c=3]{l} Mahasiswa mampu mengimplementasikan metode statistika deskriptif dan inferensial menggunakan perangkat komputasi (Python/R). &  &  \\
 & SCPMK0902-03502 & \SetCell[c=3]{l} Mahasiswa mampu menerapkan metode probabilitas, distribusi, dan pengujian hipotesis dalam analisis data TI. &  &  \\
\SetCell[r=2]{l} Penilaian & \SetCell[c=4]{l} *Nilai CPMK didapat dari agregasi nilai SUB-CPMK yang dibebankan &  &  &  \\
 & \SetCell[c=4]{l}{\tiny\begin{tblr}{width=\linewidth,colspec={|Q[c,m,0.1\linewidth]|Q[c,m,0.16\linewidth]|X[l,m]|X[l,m]|X[l,m]|X[l,m]|X[l,m]|X[l,m]|Q[c,m,0.08\linewidth]|},hlines,vlines,hspan=minimal,rowsep=1pt,colsep=0.7pt,row{1,2}={bg=headgray,font=\bfseries}}
Kode CPMK & Kode SCPMK & \SetCell[c=6]{c} Bobot per Bentuk Penilaian & & & & & & Total Bobot \\
 & & Tugas & Kuis 1 & UTS & Kuis 2 & PBL & UAS & \\
CPMK0707 & SCPMK0707-03501 & 18\%: implementasi statistika deskriptif dan inferensial & 0\% & 7\%: probabilitas dan inferensi & 5\%: regresi dan simulasi & 22\%: analisis dataset TI & 13\%: presentasi hasil analisis & 65\% \\
CPMK0902 & SCPMK0902-03502 & 12\%: probabilitas dan pengujian hipotesis & 5\%: statistika deskriptif dan probabilitas & 8\%: probabilitas dan inferensi & 0\% & 0\% & 10\%: presentasi hasil analisis & 35\% \\
\SetCell[c=8]{r}\textbf{Total Per Penilaian} & & & & & & & & \textbf{100\%} \\
\end{tblr}} &  &  &  \\
Diskripsi Singkat Mata Kuliah & \SetCell[c=4]{l} Statistik Komputasi membangun kemampuan mengaplikasikan metode statistika deskriptif, inferensial, regresi, dan simulasi menggunakan Python/R untuk analisis data dan pengambilan keputusan berbasis bukti. &  &  &  \\
Bahan Kajian / Materi Pembelajaran & \SetCell[c=4]{l}\begin{enumerate}\item Statistika deskriptif dan probabilitas\item Distribusi dan inferensi statistik\item Regresi, ANOVA, dan uji non-parametrik\item Simulasi Monte Carlo dan analisis data eksploratori\item Visualisasi data dan time series dasar\end{enumerate} &  &  &  \\
\SetCell[r=2]{l} Pustaka & \SetCell[c=4]{l} \textbf{Utama:}\begin{enumerate}\item Walpole, R. E. 2016. Probability \& Statistics for Engineers and Scientists (9th ed.). Pearson.\item VanderPlas, J. 2016. Python Data Science Handbook. O'Reilly.\item Montgomery, D. 2017. Applied Statistics and Probability for Engineers. Wiley.\end{enumerate} &  &  &  \\
 & \SetCell[c=4]{l} \textbf{Pendukung:} Materi LMS, dokumentasi resmi Python/R, dan jobsheet praktikum. &  &  &  \\
Media Pembelajaran & \SetCell[c=2]{l} \textbf{Software:} Python, R, Jupyter Notebook, RStudio, LMS. &  & \SetCell[c=2]{l} \textbf{Hardware:} Komputer/laptop, proyektor. &  \\
Nama Dosen Pengampu & \SetCell[c=4]{l} Tim Pengajar Program Studi D4 TI &  &  &  \\
Matakuliah Syarat & \SetCell[c=4]{l} - &  &  &  \\
\end{longtblr}
\begin{longtblr}{width=\textwidth,colspec={|Q[c,m,0.06\textwidth]|X[1.35,l,m]|X[1.08,l,m]|X[1.16,l,m]|X[0.7,l,m]|X[1.24,l,m]|X[1.06,l,m]|X[0.98,l,m]|Q[c,m,0.05\textwidth]|},hlines,vlines,hspan=minimal,rowhead=2,row{1,2}={bg=headgreen,font=\bfseries},rowsep=1pt,colsep=1.5pt}
Mgg.\par Ke & Kemampuan Akhir Yang Direncanakan (Sub-CP-MK) & Bahan kajian (Materi Pembelajaran) & Bentuk dan Metode Pembelajaran & Estimasi Waktu & Pengalaman Belajar Mahasiswa & Kriteria \& Bentuk Penilaian & Indikator Penilaian & Bobot Penilaian (\%) \\
(1) & (2) & (3) & (4) & (5) & (6) & (7) & (8) & (9) \\
1 & SCPMK0902-03502 & Pengantar statistika komputasi dan Python/R. & Modalitas: Blended Learning. Bentuk: Luring/Daring. Metode: Case Method, praktik komputasi, studi kasus, dan peer review. & 1 x 2 x 50' tatap muka; 1 x 2 x 50' tugas/praktik mandiri. & Mahasiswa mengerjakan aktivitas terarah untuk pengantar statistika komputasi dan python/r dan menunjukkan evidence hasil belajar. & Bentuk penilaian: Pengantar Statistika Komputasi dan Python/R. Mengacu rubrik RTM. & Ketepatan konsep; kualitas implementasi; validasi dan dokumentasi. & 2.5 \\
2 & SCPMK0902-03502 & Statistika deskriptif: mean, median, modus, variansi. & Modalitas: Blended Learning. Bentuk: Luring/Daring. Metode: Case Method, praktik komputasi, studi kasus, dan peer review. & 1 x 2 x 50' tatap muka; 1 x 2 x 50' tugas/praktik mandiri. & Mahasiswa mengerjakan aktivitas terarah untuk statistika deskriptif dan menunjukkan evidence hasil belajar. & Bentuk penilaian: Statistika Deskriptif. Mengacu rubrik RTM. & Ketepatan konsep; kualitas implementasi; validasi dan dokumentasi. & 2.5 \\
3 & SCPMK0902-03502 & Probabilitas dan distribusi diskrit. & Modalitas: Blended Learning. Bentuk: Luring/Daring. Metode: Case Method, praktik komputasi, studi kasus, dan peer review. & 1 x 2 x 50' tatap muka; 1 x 2 x 50' tugas/praktik mandiri. & Mahasiswa mengerjakan aktivitas terarah untuk probabilitas dan distribusi diskrit dan menunjukkan evidence hasil belajar. & Bentuk penilaian: Probabilitas dan Distribusi Diskrit. Mengacu rubrik RTM. & Ketepatan konsep; kualitas implementasi; validasi dan dokumentasi. & 2.5 \\
4 & SCPMK0902-03502 & Distribusi kontinu: normal, eksponensial, gamma. & Modalitas: Blended Learning. Bentuk: Luring/Daring. Metode: Case Method, praktik komputasi, studi kasus, dan peer review. & 1 x 2 x 50' tatap muka; 1 x 2 x 50' tugas/praktik mandiri. & Mahasiswa mengerjakan aktivitas terarah untuk distribusi kontinu dan menunjukkan evidence hasil belajar. & Bentuk penilaian: Distribusi Kontinu. Mengacu rubrik RTM. & Ketepatan konsep; kualitas implementasi; validasi dan dokumentasi. & 2.5 \\
5 & SCPMK0902-03502 & Kuis 1: statistika deskriptif dan probabilitas. & Modalitas: Blended Learning. Bentuk: Luring/Daring. Metode: Case Method, praktik komputasi, studi kasus, dan peer review. & 1 x 2 x 50' tatap muka; 1 x 2 x 50' tugas/praktik mandiri. & Mahasiswa mengerjakan aktivitas terarah untuk kuis 1 statistika deskriptif dan probabilitas dan menunjukkan evidence hasil belajar. & Bentuk penilaian: Kuis 1: Statistika Deskriptif dan Probabilitas. Mengacu rubrik RTM. & Ketepatan konsep; kualitas implementasi; validasi dan dokumentasi. & 5 \\
6 & SCPMK0707-03501 & Estimasi titik dan interval kepercayaan. & Modalitas: Blended Learning. Bentuk: Luring/Daring. Metode: Case Method, praktik komputasi, studi kasus, dan peer review. & 1 x 2 x 50' tatap muka; 1 x 2 x 50' tugas/praktik mandiri. & Mahasiswa mengerjakan aktivitas terarah untuk estimasi titik dan interval kepercayaan dan menunjukkan evidence hasil belajar. & Bentuk penilaian: Estimasi Titik dan Interval Kepercayaan. Mengacu rubrik RTM. & Ketepatan konsep; kualitas implementasi; validasi dan dokumentasi. & 2.5 \\
7 & SCPMK0707-03501 & Pengujian hipotesis: uji t dan uji z. & Modalitas: Blended Learning. Bentuk: Luring/Daring. Metode: Case Method, praktik komputasi, studi kasus, dan peer review. & 1 x 2 x 50' tatap muka; 1 x 2 x 50' tugas/praktik mandiri. & Mahasiswa mengerjakan aktivitas terarah untuk pengujian hipotesis uji t dan uji z dan menunjukkan evidence hasil belajar. & Bentuk penilaian: Pengujian Hipotesis: Uji t dan Uji z. Mengacu rubrik RTM. & Ketepatan konsep; kualitas implementasi; validasi dan dokumentasi. & 2.5 \\
8 & SCPMK0902-03502, SCPMK0707-03501 & UTS: probabilitas dan inferensi statistik. & Modalitas: Blended Learning. Bentuk: Luring/Daring. Metode: Case Method, praktik komputasi, studi kasus, dan peer review. & 1 x 2 x 50' tatap muka; 1 x 2 x 50' tugas/praktik mandiri. & Mahasiswa mengerjakan aktivitas terarah untuk uts probabilitas dan inferensi statistik dan menunjukkan evidence hasil belajar. & Bentuk penilaian: UTS: Probabilitas dan Inferensi Statistik. Mengacu rubrik RTM. & Ketepatan konsep; kualitas implementasi; validasi dan dokumentasi. & 15 \\
9 & SCPMK0707-03501 & Analisis regresi linear sederhana dan berganda. & Modalitas: Blended Learning. Bentuk: Luring/Daring. Metode: Case Method, praktik komputasi, studi kasus, dan peer review. & 1 x 2 x 50' tatap muka; 1 x 2 x 50' tugas/praktik mandiri. & Mahasiswa mengerjakan aktivitas terarah untuk analisis regresi linear sederhana dan berganda dan menunjukkan evidence hasil belajar. & Bentuk penilaian: Analisis Regresi Linear Sederhana dan Berganda. Mengacu rubrik RTM. & Ketepatan konsep; kualitas implementasi; validasi dan dokumentasi. & 2.5 \\
10 & SCPMK0707-03501 & ANOVA dan uji non-parametrik. & Modalitas: Blended Learning. Bentuk: Luring/Daring. Metode: Case Method, praktik komputasi, studi kasus, dan peer review. & 1 x 2 x 50' tatap muka; 1 x 2 x 50' tugas/praktik mandiri. & Mahasiswa mengerjakan aktivitas terarah untuk anova dan uji non-parametrik dan menunjukkan evidence hasil belajar. & Bentuk penilaian: ANOVA dan Uji Non-Parametrik. Mengacu rubrik RTM. & Ketepatan konsep; kualitas implementasi; validasi dan dokumentasi. & 2.5 \\
11 & SCPMK0707-03501 & Simulasi Monte Carlo. & Modalitas: Blended Learning. Bentuk: Luring/Daring. Metode: Case Method, praktik komputasi, studi kasus, dan peer review. & 1 x 2 x 50' tatap muka; 1 x 2 x 50' tugas/praktik mandiri. & Mahasiswa mengerjakan aktivitas terarah untuk simulasi monte carlo dan menunjukkan evidence hasil belajar. & Bentuk penilaian: Simulasi Monte Carlo. Mengacu rubrik RTM. & Ketepatan konsep; kualitas implementasi; validasi dan dokumentasi. & 2.5 \\
12 & SCPMK0707-03501 & Kuis 2: regresi dan simulasi. & Modalitas: Blended Learning. Bentuk: Luring/Daring. Metode: Case Method, praktik komputasi, studi kasus, dan peer review. & 1 x 2 x 50' tatap muka; 1 x 2 x 50' tugas/praktik mandiri. & Mahasiswa mengerjakan aktivitas terarah untuk kuis 2 regresi dan simulasi dan menunjukkan evidence hasil belajar. & Bentuk penilaian: Kuis 2: Regresi dan Simulasi. Mengacu rubrik RTM. & Ketepatan konsep; kualitas implementasi; validasi dan dokumentasi. & 5 \\
13 & SCPMK0707-03501 & Analisis data eksploratori (EDA). & Modalitas: Blended Learning. Bentuk: Luring/Daring. Metode: Case Method, praktik komputasi, studi kasus, dan peer review. & 1 x 2 x 50' tatap muka; 1 x 2 x 50' tugas/praktik mandiri. & Mahasiswa mengerjakan aktivitas terarah untuk analisis data eksploratori dan menunjukkan evidence hasil belajar. & Bentuk penilaian: Analisis Data Eksploratori (EDA). Mengacu rubrik RTM. & Ketepatan konsep; kualitas implementasi; validasi dan dokumentasi. & 2.5 \\
14 & SCPMK0707-03501 & Visualisasi data statistik. & Modalitas: Blended Learning. Bentuk: Luring/Daring. Metode: Case Method, praktik komputasi, studi kasus, dan peer review. & 1 x 2 x 50' tatap muka; 1 x 2 x 50' tugas/praktik mandiri. & Mahasiswa mengerjakan aktivitas terarah untuk visualisasi data statistik dan menunjukkan evidence hasil belajar. & Bentuk penilaian: Visualisasi Data Statistik. Mengacu rubrik RTM. & Ketepatan konsep; kualitas implementasi; validasi dan dokumentasi. & 2.5 \\
15 & SCPMK0707-03501 & Time series analysis dasar. & Modalitas: Blended Learning. Bentuk: Luring/Daring. Metode: Case Method, praktik komputasi, studi kasus, dan peer review. & 1 x 2 x 50' tatap muka; 1 x 2 x 50' tugas/praktik mandiri. & Mahasiswa mengerjakan aktivitas terarah untuk time series analysis dasar dan menunjukkan evidence hasil belajar. & Bentuk penilaian: Time Series Analysis Dasar. Mengacu rubrik RTM. & Ketepatan konsep; kualitas implementasi; validasi dan dokumentasi. & 2.5 \\
16 & SCPMK0707-03501 & PBL: analisis dataset TI dengan Python/R. & Modalitas: Blended Learning. Bentuk: Luring/Daring. Metode: Case Method, praktik komputasi, studi kasus, dan peer review. & 1 x 2 x 50' tatap muka; 1 x 2 x 50' tugas/praktik mandiri. & Mahasiswa mengerjakan aktivitas terarah untuk pbl analisis dataset ti dengan python/r dan menunjukkan evidence hasil belajar. & Bentuk penilaian: PBL: Analisis Dataset TI. Mengacu rubrik RTM. & Ketepatan konsep; kualitas implementasi; validasi dan dokumentasi. & 22 \\
17 & SCPMK0707-03501, SCPMK0902-03502 & UAS: presentasi hasil analisis statistik. & Modalitas: Blended Learning. Bentuk: Luring/Daring. Metode: Case Method, praktik komputasi, studi kasus, dan peer review. & 1 x 2 x 50' tatap muka; 1 x 2 x 50' tugas/praktik mandiri. & Mahasiswa mengerjakan aktivitas terarah untuk uas presentasi hasil analisis statistik dan menunjukkan evidence hasil belajar. & Bentuk penilaian: UAS: Presentasi Hasil Analisis Statistik. Mengacu rubrik RTM. & Ketepatan konsep; kualitas implementasi; validasi dan dokumentasi. & 23 \\
\end{longtblr}
